\subsection{非Ore型族の一般点は最小周期2である}\label{subsec:nonore-period}
式\eqref{eq:nonore-five}について，この小節だけ
$(K,V,B,S,A)=(k_0,v_0,b_0,s_0,a_0)$と略記する．
点対応の双線型化を用いて，接方向だけでなく実際の
非1周期的な代数を得る．

\begin{theorem}[非Ore型五変数族の一般最小周期2]\label{thm:nonore-generic-two}
式\eqref{eq:nonore-five}の基底$\mathbb A^5$には非空なZariski開集合
$\mathcal U$があり，その上で最小正周期はちょうど$2$である．
したがって\emph{どのgraded algebraにも由来しない}二次元の
同型類の代数的族が得られる．さらにこの族の一般点の左点スキームは
一次元の支持を持つため，十二個の単純点と一個の非被約点を持つ
九変数楕円Clifford--Ore族の一般点とは非同型である．
ただし二つの族の閉包が異なる既約成分であるとは主張しない．
\end{theorem}
\begin{proof}
位置$i$で$\varepsilon=(-1)^i$とする．左点スキーム$E_\varepsilon$は，
始点座標$[w:y:z:x]\in\mathbb P^3$に対し，次の$6\times4$行列の
階数が$3$以下となる零点スキームである．
\begin{equation}\label{eq:nonore-point-matrix}
 M_\varepsilon=\begin{pmatrix}
 y&w&0&-\varepsilon Az\\
 z&0&w&0\\
 -Kw&z&y&0\\
 x&0&-Vz&w\\
 -Bw&x&-Sz&y\\
 0&0&x&z
 \end{pmatrix}.
\end{equation}
$AKVBS\ne0$とし，次の五次式は重根を持たないと仮定する．
\begin{equation}\label{eq:nonore-quintic}
 Q_\varepsilon(z)=\varepsilon AV^2z^5+4Sz^2-4B.
\end{equation}
全$4\times4$小行列式を調べると，$E_\varepsilon$の支持は一本の直線
\[
 \ell=\{w=z=0\}
\]
および，$Q_\varepsilon(z)=0$を満たす相異なる五つの点
\begin{equation}\label{eq:nonore-five-points}
 p_\varepsilon(z)=
 \left[1:-\frac K{2z}-\frac{\varepsilon AV}4z^3:
              z:-\frac V2z^2\right]
\end{equation}
に正確に一致する．必要性を得るには，例えば行列の行番号を$0$から数え，
$(1,2,3,5)$行の小行列式
$-z^2(Vz^2+2wx)$，$(0,1,2,5)$行の
$z(\varepsilon Axz^2-Kw^3-2wyz)$，$(0,1,4,5)$行の
$z(\varepsilon Ax^2z-Bw^3+Swz^2)$を用いればよい．
$z=0$なら$(0,1,2,3)$行の小行列式から$Kw^4=0$，従って$w=0$となる．
$z\ne0$では第一式により$w\ne0$なので$w=1$とでき，
後二式が\eqref{eq:nonore-five-points}と\eqref{eq:nonore-quintic}を与える．
逆に式\eqref{eq:nonore-five-points}を代入した全15小行列式は
$Q_\varepsilon$で割り切れる．$w=z=0$上でも全15小行列式は零になる．
これらの恒等式は\path{computations/verify_nonore_point_geometry.py}で，
パラメータを特殊化せず厳密に確認する．
従って$\ell$は唯一の正次元既約支持成分であり，
任意の$E_+\simeq E_-$の射影線型同型は$\ell$を固定する．

$Q_-( -z)=Q_+(z)$なので，$E_-$の五つの孤立点の
直線$\ell$からの射影座標は，$E_+$の五つの根の負である．
ここで一般の$Q_+$の根の集合には$\operatorname{PGL}_2$の非自明な
自己同型がないことを確認する．非空性は
\[
 K=V=A=1,\qquad B=-\tfrac14,\quad S=\tfrac14
\]
という有理数点で示せる．この点では$Q_+(z)=z^5+z^2+1$である．
この多項式は$\mathbb F_2$上既約で，$\mathbb F_5$上では
\[
 (z^2-z+2)(z^3+z^2-z-2)
\]
と，次数$(2,3)$の相異なる既約因子に分かれる．従ってガロア群は
推移的な5次置換群で，5-cycleと$(2,3)$型の置換を含み，
特にtranspositionを含む．推移的な5次置換群は素数次数ゆえ原始的であり，
原始置換群にtranspositionがあれば，その共役の作る辺グラフは連結で，
対応するtranspositionが$S_5$を生成する．従ってガロア群は$S_5$である．

五根の$\operatorname{PGL}_2$安定化群$\Gamma$は，五根への作用として
$S_5$へ忠実に埋まる．ガロア共役で$\Gamma$は保たれるから，
その像は$S_5$の正規部分群である．非自明なら$A_5$を含むが，
その五点作用で一点を固定する部分群は$A_4$を含む．
これは有限な$\operatorname{PGL}_2$部分群の一点安定化群が，
標数零では巡回群でなければならないことに矛盾する．
よって$\Gamma=1$である．異なる五点に対する非自明な置換が
$\operatorname{PGL}_2$で実現する条件は，各置換ごとの閉条件である．
五点の有限\'etale被覆から降下すれば，安定化群が自明な領域は
非空なZariski開集合となる．以下この開集合上で議論する．

$\ell$を固定する射影変換は商平面の射影直線
$[w:z]$にM\"obius変換を誘導する．
五根の集合をその負の集合へ送るM\"obius変換は，
上の安定化群の自明性から$z\mapsto-z$ただ一つである．
従ってもし$E_+\simeq E_-$なら，
$w,z$の二座標上では$[w:z]\mapsto[w:-z]$であり，
その$y$座標は五つの根上で
$1,z,z^2,y_+(z)$の線型結合として$y_-(-z)$に一致しなければならない．

ところが$Q_+(z)=0$より
\[
 z^{-1}=\frac{AV^2z^4+4Sz}{4B}
\]
である．従って$k[z]/(Q_+)$の部分空間
$\langle1,z,z^2\rangle$で割ると，
\[
 \begin{aligned}
 y_+(z)&\equiv-\frac{AV}{4}z^3-\frac{KAV^2}{8B}z^4,\\
 y_-(-z)&\equiv-\frac{AV}{4}z^3+\frac{KAV^2}{8B}z^4.
 \end{aligned}
\]
この二つの係数ベクトルの行列式は
\[
 -\frac{KA^2V^3}{16B}\ne0.
\]
すなわち後者は$1,z,z^2,y_+(z)$の線型包に入らない．
相異なる五根での関数の一致は$k[z]/(Q_+)$での一致と同値なので，
上で必要とされた射影変換は存在しない．
従って$E_+$と$E_-$は射影同型でない．

1周期同型があれば，隣接する位置の左点スキームが射影同型になるから，
これは1周期性を排除する．全ファイバーには厳密2周期があるので，
最小正周期は2である．同型類の像が二次元であることは
定理\ref{thm:nonore-five}ですでに証明した．
最後に，一般の九変数楕円Clifford--Ore族の左点スキームは
非被約一点と十二個の単純点という零次元支持を持つのに対し，
本族の左点スキームは直線$\ell$を含む．
同型は点スキームを保つので，両一般点は非同型である．
\end{proof}

\begin{remark}
一般の二周期性および九変数楕円族の一般点との非同型性は上で証明した．
別の三生成部分を選ぶと本族はOre表示を持ち，七変数族への閉包包含も
定理\ref{thm:nonore-in-seven-closure}で証明される．
両族が異なる既約成分をなすという主張はしない．
\end{remark}
